\subsection{Extraneous Solutions} 
The domain of the logarithm function contains only \makeblank[1in] numbers. If we multiply or divide two \makeblank[1in] numbers, the result is \makeblank[1in]. So when we \makeblank[1in] logarithms, we risk introducing extraneous solutions.
\begin{Example}
Solve $\log_4(x+7)=3-\log_4(x-5)$. Discard any extraneous solutions.
\begin{lpic}{}{8}
\regtextleft{1}{-1}{Collect the log terms:};
\regtextleft{1}{-3}{Combine the logs:};
\regtextleft{1}{-4}{Rewrite in exponential form:};
\regtextleft{1}{-5}{Finish solving:};
\end{lpic}
To identify extraneous solutions, we plug in:
\begin{lpic}{}{4}
\end{lpic}
Notice that the extraneous solution was not in the \makeblank[1in] of one of the log functions. We could also note the restrictions on $x$ up front:
\begin{lpic}{}{2}
\end{lpic}
We could have gone in knowing that only solutions such that \makeblank[1in] would be true solutions.
\end{Example}